% abstract.tex
\begin{abstract}
\noindent
Personal financial management depends on the dull, error-prone discipline of
recording every transaction. Most people abandon expense-tracking applications
not because the analytics are poor but because the \emph{data-entry friction} is
too high: navigating forms, dropdowns, and keyboards for each coffee or taxi
fare is incompatible with the moments in which spending actually happens. This
project, \textbf{MiPay}, attacks that friction directly with a
\emph{voice-first} interaction model: the user presses a single button, speaks a
natural sentence describing money spent or received---in Arabic, English, or a
code-switched mixture of both---and the system returns a structured, editable
transaction ready to save.

The central academic contribution is that the entire artificial-intelligence
pipeline is \emph{self-hosted and built from open-source models}, rather than a
thin wrapper over a commercial cloud API. Speech is transcribed by a locally
hosted Whisper model served through \texttt{faster-whisper} (a CTranslate2
re-implementation that is roughly four times faster than the reference
implementation on CPU). The transcript is then converted into structured
financial fields---transaction type, amount, currency, category, counterparty,
and date---by a 7-billion-parameter open-weight instruction-tuned large language
model (Qwen2.5-7B-Instruct) served locally through Ollama. Crucially, the model
is invoked with \emph{JSON-schema-constrained decoding}, which guarantees that
every response is syntactically valid and conforms to the target schema,
eliminating an entire class of parsing failures common to prompt-only JSON
generation. A deterministic, unit-tested post-processing layer then handles the
tasks that language models perform unreliably---Arabic-Indic numeral
normalization, relative-date resolution (``yesterday'', the Arabic ``ams'',
``last Friday''), currency defaulting, and a confidence-aware
\texttt{needs\_review} decision---before a human-in-the-loop confirmation step
commits anything to the database.

The system is delivered as a Flutter mobile application backed by a FastAPI
service, PostgreSQL, and Ollama, fully containerized with Docker Compose, and
supports three input modes (voice, typed text, and fully manual entry) that all
converge on the same confirmation and save path. The self-hosted design yields
three properties that a cloud wrapper cannot claim: \emph{zero per-request
inference cost}, \emph{full data privacy} (audio is deleted immediately after
transcription and never leaves the operator's machine), and a genuine
\emph{model-selection and evaluation methodology} suitable for academic scrutiny.

We evaluate the pipeline on a purpose-built corpus of more than 260 labeled
bilingual utterances using Word Error Rate for transcription and per-field
exact-match accuracy and F\textsubscript{1} for extraction, under two conditions
(gold-transcript and full end-to-end audio), together with ablations across
Whisper model sizes, language-model sizes, and few-shot versus zero-shot
prompting. The thesis documents the system as built and articulates a concrete
upgrade path---larger and newer speech and language models, and task-specific
fine-tuning on user-corrected data---as principled future work.

\vspace{0.4cm}
\noindent\textbf{Keywords:} speech recognition, large language models,
information extraction, constrained decoding, Arabic natural language processing,
code-switching, personal finance, edge AI, privacy-preserving machine learning,
Flutter, FastAPI.
\end{abstract}
